\paragraph{Finite program capacities.}
The omission bound in Theorem~\ref{thm:omission} extends to programs with integer capacities $u_j$ and strict responsive priorities. Let $H$ again be the applicants whose original assignment is omitted from their display, and $K_j,T_j$ their numbers at program $j$ before and after the restriction. Both gross placement losses and gains are bounded by
\[
                         |H|-\sum_j\min\{K_j,T_j\}.
\]
To see this, removing the members of $H$ and their occupied seats leaves a stable assignment of the other applicants under the residual capacity vectors $u-K$ and $u-T$, respectively. Any block in the residual would have blocked the original. These applicants have the same fixed displayed lists in both residual markets; their original surviving assignments on those lists remain stable after restriction. Raising a single program's capacity by one adds at most one newly matched applicant and removes none. One can replace each capacity by copies with the program's priority, rank copies together in a common fixed order, and insert the last copy by program-proposing DA as in Theorem~\ref{thm:occupied-seats}. Matched-set invariance transfers that statement across stable selections. Comparing the two residual stocks through their coordinatewise minimum bounds incumbent losses by $\sum_j(T_j-K_j)_+$ and gains by $\sum_j(K_j-T_j)_+$. All members of $H$ originally held a place, and $|H|-\sum_jT_j$ lose one, giving the two bounds. Crucially, a recommendation exposes a program with all its seats at once. The quota is a count of programs and the overlap term is a count of seats. These assumptions exclude dynamic sibling priorities and type-specific reserve choice functions.

The protection certificate in Theorem~\ref{thm:threat} also has a many-seat implementation. Build a right vertex for \emph{each protected baseline holder}, keeping their program label. A new outside applicant has an edge to a holder if that applicant sees the holder's program and could outrank the holder there. At most the matching number of this graph of persons can lose their placement. If there are no such edges, every protected holder weakly prefers their new applicant-proposing assignment to their baseline program. To prove both statements, clone program $j$ into $u_j$ seats. All agents rank its seats consecutively in a common order, each with the program's priority. Lift any stable assignment by placing its priority-ordered holders in the consecutively preferred seats. A same-program blocking pair is precluded by this ordering and any cross-program blocking pair would have blocked the original assignment. Conversely, any clone-stable matching induces a stable program assignment. Applicants see all seats of a displayed program and have unchanged rankings across programs. Applying the corresponding unit-seat results proves the count and rank claims; applicant-optimality on the cloned market projects to applicant-optimality over programs. With unknown priorities, aggregate the holders at a program into a single vertex of capacity equal to their number. With known strict priorities the potential targets at one program form a suffix of its holder list. A flow network can share suffix arcs rather than duplicating one edge for every possible holder.
